In this paper, we tried to bring a comprehensive framework for detecting 3D chessboards and identifying individual chess pieces in real-world chess games.
The approach proposed combines traditional computer vision techniques with deep learning methods such as YOLOv8m to address the multiple challenges of variable lighting,
perspective distortion and heterogeneous piece design recognition. \newline
At the end of our study, we selected a top-down contour-based strategy that demonstrated superior robustness under real-world constraints.
About the object detection part, we conducted extensive experimentation with both transformed-based and convolutional architectures, determining that YOLOv8m, trained on a large dataset, offers the most accurate and reliable performance for the task of chess-piece classification.
\newline
After all the recognition part, we developed a practical pipeline for generating FEN (Forsyth-Edwards Notation) notation from a board state, enabling seamless conversion from a physical board states to a standard digital representation.
In addition, we proposed a binary embedding-based similarity search method for retrieving similar state games  from a large dataset of millions of games.
\newline
Despite the promising results, some limitation remains.
Particularly the reliance on a full board visibility.
